\section{Study Design Explained}

\subsection{Open-label, single-arm, and what that means for you}

Three technical terms describe the design, and each has a consequence a
participant can act on.

\textbf{Open-label} means nothing about the treatment is concealed. The
participant knows the dose level they are on, the software version that will
operate, and the identity of the operating surgeon. There is no masking to
maintain and therefore nothing that has to be withheld to protect the study.

\textbf{Single-arm} means there is one treatment group and no allocation. The
concern reported most often about trial mechanics is being assigned to a
treatment the participant would not have chosen \cite{NCITrial2024,
WuSemiAI2026}. In this study that cannot happen: every participant who enrols
receives the same intervention, and the only thing that varies between them is
the daraxonrasib dose level, which is set by what happened to the people ahead
of them rather than by chance.

\textbf{Dose-finding} means the study is looking for the highest dose that is
tolerated, not for the largest effect. A participant at DL1 is not receiving a
lesser treatment; they are receiving the dose that the study is confident enough
to start with, and their tolerance of it is what permits DL2 to open.

\subsection{Your status in the study, and how it can change}

Figure 12 draws participant status as a state machine, and the reason to use a
state machine rather than a flowchart is the guard. Every transition in the
figure carries a named condition that must evaluate true before the transition
can be taken, and the protections live in the guards rather than in the states.
A participant who wants to know what stands between them and the operating room
should read the guard on the transition into the operative state:
\texttt{[Phase 0 gate signed]}, \texttt{[USL $\geq$ 7.0 verified]}, and
\texttt{[pre-procedure matrix clear]}, all three, conjunctively.

One state in that figure exists only for the participant's benefit. The cohort
hold is entered when a prior participant is still inside their sentinel window,
and it costs the study calendar time. It would be tidier to fold it into a
transition guard, and it is drawn as an explicit state because it is the
clearest single piece of evidence that this design accepts delay in exchange for
safety.

\begin{pafloat}
\begin{pafig}[plantuml-type: state machine, @startuml]
\node[umlinit] (i0) at (0,1.5) {};
\node[umlstate,text width=44mm] (s1) at (0,0)     {\textbf{Screened}\\screening, day $-28$ to $-1$; KRAS G12 confirmed, staging done};
\node[umlstategray,text width=36mm] (sx) at (-7.4,0) {\textbf{Not enrolled}\\screen failure or declined; minimal data set only, \S 5.4; standard care unchanged};
\node[umlstatesoft,text width=44mm] (s2) at (0,-2.9) {\textbf{Enrolled}\\consent on file; software version frozen and recorded};
\node[umlstategray,text width=36mm] (sh) at (7.6,-2.9) {\textbf{Cohort hold}\\sentinel window open; no further participant dosed until it closes};
\node[umlstatesoft,text width=44mm] (s3) at (0,-5.9) {\textbf{Dose assigned}\\DL1, DL2, or DL3 by the 3+3 rule; never above a level not yet cleared};
\node[umlstateon,text width=44mm] (s4) at (0,-8.9)  {\textbf{Operative}\\day 0, robotic Whipple under continuous human oversight};
\node[umlstateon,text width=44mm] (s5) at (0,-11.6) {\textbf{Recovering}\\days 1 to 30; daraxonrasib restart advisory at T+7, T+14, or T+21};
\node[umlstatesoft,text width=44mm] (s6) at (0,-14.3){\textbf{Confirmed}\\day 90; R0 pathology and 90-day mortality};
\node[umlstate,text width=44mm] (s7) at (0,-17.0)  {\textbf{Long term}\\every 12 weeks to 24 months; PFS and OS};
\node[umlfinal] (f0) at (0,-18.6) {};
\node[umlfinal] (fx) at (-7.4,-1.7) {};
\draw[umlarrow] (i0) -- (s1);
\draw[umlarrow] (s1) -- node[umlguard,above]{[ineligible OR declined]} (sx);
\draw[umlarrow] (sx) -- (fx);
\draw[umlarrow] (s1) -- node[umlguard,pos=0.5,align=left,xshift=32mm]{[all \S 5.1 met]\\{}[no \S 5.2 exclusion]\\{}[consent signed]} (s2);
\draw[umlarrow] (s2) to[out=0,in=90,looseness=0.9] node[umlguard,pos=0.55,align=left]{[prior participant still inside the sentinel window]} (sh);
\draw[umlarrow] (sh) to[out=-90,in=0,looseness=0.9] node[umlguard,pos=0.30,align=left]{[sentinel window closed]\\{}[no DLT declared]} (s3);
\draw[umlarrow] (s2) -- node[umlguard,pos=0.5,align=left,xshift=30mm]{[no sentinel window open]} (s3);
\draw[umlarrow] (s3.west) to[out=170,in=190,looseness=5] node[umlguard,left,align=right]{[DLT observed in cohort]\\/ expand 3 to 6, do not escalate} (s3.south west);
\draw[umlarrow] (s3) -- node[umlguard,pos=0.5,align=left,xshift=32mm]{[Phase 0 gate signed]\\{}[USL $\geq 7.0$ verified]\\{}[pre-procedure matrix clear]} (s4);
\draw[umlarrow] (s4) -- node[umlguard,pos=0.5,align=left,xshift=30mm]{[procedure complete OR converted to open]} (s5);
\draw[umlarrow] (s5) -- node[umlguard,pos=0.5,align=left,xshift=26mm]{[day-30 window closed]} (s6);
\draw[umlarrow] (s6) -- node[umlguard,pos=0.5,align=left,xshift=26mm]{[day-90 assessments done]} (s7);
\draw[umlarrow] (s7) -- node[umlguard,pos=0.5,align=left,xshift=30mm]{[24-month endpoint OR withdrawal OR death]} (f0);
\node[umlnote,text width=42mm,anchor=north west] at (5.4,-6.5)
  {This state exists so that you are never the second person at a new dose level before the first person's safety window has closed. It costs the study time. It is not there for the study's benefit.};
\node[umlnote,text width=44mm,anchor=north east] at (-5.0,-8.2)
  {Single arm, open label. There is no allocation that could assign you to a comparator, and nothing about your treatment is concealed from you.};
\node[umlnote,text width=44mm,anchor=north east] at (-5.0,-12.4)
  {Conversion to open surgery is a normal outcome of this state, not a failure of it. It is recorded, reported, and does not end your participation.};
\end{pafig}
\vspace{-0.7cm}
\figcaption{Figure 12. Participant status as a guarded state machine, screening through the 24-month\\
endpoint, with the sentinel cohort hold drawn as an explicit state, because it is the\\
clearest evidence that the design accepts calendar delay in exchange for safety.}
\end{pafloat}

\subsection{The 3+3 rule, and why you are never the untested dose}

The rule is usually stated statistically and is easier to check operationally.
Three participants are treated at a dose level. If none has a dose-limiting
toxicity, the next level opens. If one does, the level is expanded to six rather
than escalated; if that sixth participant brings the total to one of six, the
next level opens, and if two or more of six have a toxicity, escalation stops
and the level below becomes the recommended dose. If two or three of the first
three have a toxicity, escalation stops immediately.

Layered on top is the sentinel stagger. The second participant at a new dose
level is not dosed until the first participant's safety window has closed. That
delay is the difference between a rule that protects the cohort and a rule that
protects the individual, and it is why the answer to "am I the untested dose" is
no: the dose a participant receives was cleared by the people ahead of them, and
the study expands rather than escalates the moment anyone is harmed.

Figure 13 draws the rule as a decision tree, which is the right structure
because there is exactly one path from the root to each leaf. Two of the leaves
are stops, and one of those convenes the Data and Safety Monitoring Board and is
reported to every enrolled participant. A dose-escalation figure with no
stopping leaf would misrepresent the rule.

\begin{pafloat}
\begin{pafig}[graphviz-type: digraph, decision tree with record leaves]
\node[gvkey,text width=30mm] (root) at (0,0) {Dose level 1\\first 3 participants\\sentinel staggered};
\node[gvsoft,text width=22mm] (d10) at (-5.6,-3.1) {0 of 3 had a DLT};
\node[gvgray2,text width=22mm](d11) at (0,-3.1)    {1 of 3 had a DLT};
\node[gvgray3,text width=22mm](d12) at (5.6,-3.1)  {2 or 3 of 3 had a DLT};
\node[gvsoft,text width=22mm] (ex)  at (0,-6.0)    {expand this level to 6 participants};
\node[gvsoft,text width=22mm] (e1)  at (-2.9,-8.8) {1 of 6 total};
\node[gvgray3,text width=22mm](e2)  at (2.9,-8.8)  {2 or more of 6};
\node[gvmid,text width=24mm]  (d2)  at (-7.6,-6.6) {Dose level 2\\next 3, sentinel staggered};
\node[gvmid,text width=24mm]  (d3)  at (-7.6,-10.4){Dose level 3\\next 3, sentinel staggered};
\node[gvbox,text width=44mm,minimum height=13mm] (stopb) at (5.2,-12.4)
  {\textbf{STOP ESCALATING}\\this level is above the tolerated dose; the level below becomes the recommended dose; no participant is dosed higher};
\node[gvbox,text width=44mm,minimum height=13mm] (mtd) at (-3.0,-14.8)
  {\textbf{RECOMMENDED PHASE 2 DOSE}\\the highest level at which fewer than 2 of 6 participants had a dose-limiting toxicity, reported with the safety data};
\node[gvboxd,text width=28mm,minimum height=10mm] (halt) at (10.0,-8.2)
  {Study pause\\DSMB convenes\\you are told};
\begin{scope}[on background layer]
\node[gvcluster,fit=(d10)(d11)(d12),inner sep=8pt] (SEN) {};
\end{scope}
\node[gvctitle,anchor=south west] at (SEN.north west) {sentinel window: nobody advances until this closes};
\draw[gvedge] (root) to[out=-150,in=90,looseness=0.8] node[umlguard,pos=0.35]{0 of 3} (d10);
\draw[gvedge] (root) -- node[umlguard,pos=0.4]{1 of 3} (d11);
\draw[gvedge] (root) to[out=-30,in=90,looseness=0.8] node[umlguard,pos=0.35]{$\geq 2$ of 3} (d12);
\draw[gvedge] (d10) to[out=-135,in=60,looseness=0.85] node[umlguard,pos=0.4]{escalate} (d2);
\draw[gvedge] (d11) -- node[umlguard,pos=0.45]{do not escalate} (ex);
\draw[gvedge] (ex) to[out=-135,in=90,looseness=0.85] node[umlguard,pos=0.4]{1 of 6} (e1);
\draw[gvedge] (ex) to[out=-45,in=90,looseness=0.85] node[umlguard,pos=0.4]{$\geq 2$ of 6} (e2);
\draw[gvedge] (e1) to[out=180,in=-20,looseness=0.85] node[umlguard,pos=0.5]{escalate} (d2);
\draw[gvedge] (e2) to[out=-45,in=150,looseness=0.85] node[umlguard,pos=0.4]{stop} (stopb);
\draw[gvedge] (d12) to[out=-60,in=120,looseness=0.85] node[umlguard,pos=0.4]{stop} (stopb);
\draw[gvedged] (d12) to[out=0,in=90,looseness=0.9] node[umlguard,pos=0.4]{if at the first level} (halt);
\draw[gvedgeb] (d2) -- node[umlguard,pos=0.5,align=left,xshift=24mm]{0 of 3, or 1 of 6} (d3);
\draw[gvedged] (d2) to[out=-40,in=170,looseness=0.9] node[umlguard,pos=0.6]{$\geq 2$ of 6} (stopb.west);
\draw[gvedgeb] (d3) to[out=-60,in=180,looseness=0.85] node[umlguard,pos=0.4]{0 of 3, or 1 of 6} (mtd);
\draw[gvedgeb] (stopb) to[out=180,in=0,looseness=0.85] node[umlguard,pos=0.45]{the level below is the answer} (mtd);
\end{pafig}
\vspace{-0.7cm}
\figcaption{Figure 13. The 3+3 dose-escalation rule drawn as a decision tree with exactly one path from\\
the root to each leaf, including the two stopping leaves and the study pause that convenes\\
the Data and Safety Monitoring Board and is reported to every participant.}
\end{pafloat}

\subsection{The dose levels}

\vspace{0.30em}

\begin{xltabular}{\textwidth}{L{1.9cm} L{3.1cm} C{1.9cm} Y}
\toprule
\textbf{Level} & \textbf{Daraxonrasib dose} & \textbf{Cohort size} & \textbf{What has to be true before this level opens} \\
\midrule
\endfirsthead
\multicolumn{4}{@{}l}{\footnotesize\itshape Table continued from the previous page.}\\
\toprule
\textbf{Level} & \textbf{Daraxonrasib dose} & \textbf{Cohort size} & \textbf{What has to be true before this level opens} \\
\midrule
\endhead
\midrule
\multicolumn{4}{r@{}}{\footnotesize\itshape Table continued on the next page.}\\
\endfoot
\bottomrule
\endlastfoot
DL1 & Starting dose level & 3, expandable to 6 & Phase 0 gate signed, USL at least 7.0, IND in effect, IDE approved by both the FDA and the IRB \\
DL2 & One step above DL1 & 3, expandable to 6 & DL1 cleared: 0 of 3, or 1 of 6, dose-limiting toxicities, and every sentinel window at DL1 closed \\
DL3 & Two steps above DL1 & 3, expandable to 6 & DL2 cleared on the same rule, plus a DSMB review at the cohort boundary \\
\end{xltabular}

\vspace{0.30em}

\noindent Planned enrollment is up to 18 participants across the three levels.
Enrollment is staggered: no participant enters a level while a sentinel window
at that level remains open. The dose anchor for the starting level is taken from
the pan-KRAS programme in previously treated metastatic disease
\cite{rasolute302, oreilly2026darax}, adjusted for the perioperative setting.

\subsection{The Phase 0 gate that runs before any patient}

No patient-facing robotic activity occurs until a simulation-validation gate has
been signed. The gate has three components, and all three must pass.

\textbf{Volume and independence.} At least 1000 simulated procedures across at
least two independent simulation frameworks. Independence matters more than
volume: two frameworks that share a physics engine share its errors, and the
protocol requires that they do not.

\textbf{Fidelity.} Motion reproduced within a 2 mm positional tolerance and a
0.5 N force tolerance across frameworks. These are the credibility questions
that the ASME verification and validation standard for computational modelling
in medical devices \cite{asmevv40} and the NASA models-and-simulations standard
\cite{nasastd7009} were written to frame.

\textbf{Readiness.} A documented Unified Safety Level of at least 7.0 for
surgical deployment, under the Physical AI overlay adapted for end-to-end
oncology investigations \cite{cfr312paiuni}. The simulation source for the
operative values is the author's 60-second robotic Whipple and daraxonrasib
simulation \cite{pdac060s2030}.

The gate sign-off is a dated document and the participant is shown it at the
baseline visit. Section 10.3 reports what happened across those 1000 runs,
including the 8.2 percent that were halted by a safety mechanism, and states in
the same sentence that a simulated worst case is not a human worst case.

\subsection{Stopping rules, stated before they are needed}

Prespecified triggers pause or stop the study. Four of them are device-specific,
and those are the ones a participant worried about the robot will look for
first. Any breach of a per-arm or cumulative force cap. Any entry into a
vascular no-fly corridor. Any emergency-stop latency measured above budget. Any
loss of the heartbeat bus that is not resolved by the automatic halt.

\vspace{0.30em}

\begin{xltabular}{\textwidth}{L{4.6cm} L{2.6cm} Y}
\toprule
\textbf{Trigger} & \textbf{Effect} & \textbf{Who decides, and who is told} \\
\midrule
\endfirsthead
\multicolumn{3}{@{}l}{\footnotesize\itshape Table continued from the previous page.}\\
\toprule
\textbf{Trigger} & \textbf{Effect} & \textbf{Who decides, and who is told} \\
\midrule
\endhead
\midrule
\multicolumn{3}{r@{}}{\footnotesize\itshape Table continued on the next page.}\\
\endfoot
\bottomrule
\endlastfoot
Two or more DLTs at a dose level & Escalation stops & Sponsor-Investigator declares; DSMB reviews; every enrolled participant is told \\
Any device-related death & Immediate study halt & DSMB; FDA and IRB notified within the reporting window \\
Force-cap breach at the tip & Procedure halt, then study pause & The gate halts automatically; the DSMB decides on resumption \\
Entry into a vascular no-fly corridor & Procedure halt, then study pause & The gate halts automatically; Physical AI Safety Review Committee re-gates \\
Emergency-stop latency above budget & Study pause pending re-gate & Physical AI Safety Review Committee; USL reassessed \\
Heartbeat loss not resolved by halt & Study pause pending re-gate & Physical AI Safety Review Committee; USL reassessed \\
Grade B/C fistula rate above the phase-1 comparator level & DSMB review at the cohort boundary & DSMB; the rate is published with the safety data \\
\end{xltabular}

\vspace{0.30em}

\noindent In every row, the participant is told. That is the point of publishing
the table before it is needed: a stopping rule that a participant only learns
about after it fires is a governance mechanism, and a stopping rule they were
shown in advance is a commitment.

\subsection{How an adverse event is reported, and how fast}

Reporting timelines are usually written for the regulator. Restated for the
participant, they answer a different question: how long after something happens
to me does somebody outside this hospital know about it.

\vspace{0.30em}

\begin{xltabular}{\textwidth}{L{4.3cm} C{2.1cm} L{3.2cm} Y}
\toprule
\textbf{Event class} & \textbf{Reported within} & \textbf{Reported to} & \textbf{What you are told, and when} \\
\midrule
\endfirsthead
\multicolumn{4}{@{}l}{\footnotesize\itshape Table continued from the previous page.}\\
\toprule
\textbf{Event class} & \textbf{Reported within} & \textbf{Reported to} & \textbf{What you are told, and when} \\
\midrule
\endhead
\midrule
\multicolumn{4}{r@{}}{\footnotesize\itshape Table continued on the next page.}\\
\endfoot
\bottomrule
\endlastfoot
Death or life-threatening event, related and unexpected & 7 calendar days & FDA, IRB, DSMB & Immediately, and in writing before discharge \\
Other serious, related, and unexpected event & 15 calendar days & FDA, IRB, DSMB & Within 24 hours, by the Site PI \\
Unanticipated adverse device effect & 10 working days & FDA, IRB, DSMB & Within 24 hours, by the Site PI \\
Any safety-limit event, force cap or no-fly & Same procedure & Physical AI Safety Review Committee & Before you leave hospital, with the measured value \\
Any emergency-stop activation & Same procedure & Sponsor-Investigator, DSMB & Before you leave hospital, with the measured latency \\
Protocol deviation affecting you & 5 working days & IRB & Within 5 working days, in writing \\
\end{xltabular}

\vspace{0.30em}

\noindent The last three rows are the device-specific ones and they have no
counterpart in a conventional drug protocol. They exist because a safety-limit
event is not an adverse event in the regulatory sense: nothing may have happened
to the participant at all. The protocol reports them anyway, and tells the
participant, because a limit that is never reported when it fires is
indistinguishable from a limit that does not exist.

\subsection{What a dose-limiting toxicity is, in this study}

The term decides whether a dose level opens, so it needs a definition a
participant can check rather than a category name. A dose-limiting toxicity in
this study is any of the following, occurring within the day-30 window and
attributed to daraxonrasib rather than to the operation.

\vspace{0.30em}

\begin{xltabular}{\textwidth}{L{5.2cm} Y}
\toprule
\textbf{Event} & \textbf{Why it counts} \\
\midrule
\endfirsthead
\multicolumn{2}{@{}l}{\footnotesize\itshape Table continued from the previous page.}\\
\toprule
\textbf{Event} & \textbf{Why it counts} \\
\midrule
\endhead
\midrule
\multicolumn{2}{r@{}}{\footnotesize\itshape Table continued on the next page.}\\
\endfoot
\bottomrule
\endlastfoot
Grade 4 haematologic toxicity lasting more than 7 days & Marrow suppression at this duration delays adjuvant therapy \\
Grade 3 or higher non-haematologic toxicity, excluding controllable nausea and fatigue & The exclusions are the toxicities that respond to supportive care \\
Any toxicity delaying the planned operation by more than 14 days & A delay of this length changes the oncologic calculus \\
Any toxicity preventing the restart of the drug by T+21 & The restart window is the endpoint the advisory layer targets \\
Grade 3 or higher hepatic transaminase elevation & The operation loads the liver; the drug must not compound it \\
\end{xltabular}

\vspace{0.30em}

\noindent Attribution between drug and operation is adjudicated by the Site
Principal Investigator with DSMB review at the cohort boundary, and the
attribution is recorded with its reasoning. That reasoning is in the
participant's own audit extract, described in \S 11.5.

\subsection{Why there is no control group, and what replaces one}

The most frequent design objection a participant raises is the absence of a
control group, and it deserves a direct answer rather than a methodological
one. Without a control group there is nothing to compare the study's results
against, which appears to make the results uninterpretable.

The answer has three parts. First, a Phase 1 study is not asking a comparative
question. It is asking whether the combination can be delivered at all and
whether it hurts people, and both of those are answerable from the single arm.
Second, randomising a participant to a conventional Whipple when the platform's
feasibility is unknown would expose the randomised group to a study's burden
without any prospect of the study's benefit, which is the reason ethics review
does not permit it at this stage. Third, and most practically, an external
comparator replaces the internal one.

That comparator is a published, contemporary, nationwide robotic
pancreaticoduodenectomy series, and it is used in one specific way: not to
declare superiority, which \S 10.5 shows is impossible at this size, but to set
the boundary beyond which the study stops. If the R0 rate falls materially below
the comparator, or the ISGPS grade B/C fistula rate rises materially above it,
the stopping rules in \S 5.6 apply. The comparator is therefore a floor rather
than a target, and every number taken from it is labelled C for comparator in
the provenance table in \S 10.6, never presented as a result of this study.

The corollary is worth stating in the participant's terms. A participant in this
study is not being compared with anyone. They are being watched against a
published floor, and if the study drops through it, the study stops before the
next person is enrolled.

\subsection{What happens if the study is stopped while you are in it}

Stopping rules are written from the sponsor's side, and read from the sponsor's
side they describe an orderly wind-down. From the participant's side the
question is narrower: what happens to me, specifically, on the day it stops.

If the study halts before the operation, the participant returns to standard
care with the study's staging imaging and tissue review in hand, and the site
makes the referral rather than leaving it to the participant. Nothing about the
halt changes the participant's eligibility for a conventional resection, and the
work-up already done does not have to be repeated.

If the study halts after the operation, the operation is complete and is not
undone by the halt. Follow-up continues, because the safety endpoints that
matter to the participant, the day-30 and day-90 assessments, are the ones that
protect them. Daraxonrasib is stopped or continued according to the reason for
the halt: a halt for a device reason does not by itself stop a drug the
participant is tolerating, and a halt for a drug reason does.

In both cases every enrolled participant is told, whether or not the halt arose
from their own case, and the notification states the reason rather than the fact
alone. That is one of the three affirmative obligations in Figure 2, and it is
not discretionary.

Data already collected remains in the study record unless the participant
withdraws it under the routes in \S 8, and \S 8.1 states which of those routes
can and cannot remove data that has already been analysed.
